\chapter{動的定式化}
\label{c-dynamic}

本章では、コストが測地線距離の二乗である場合の optimal transport の測地線的（動的とも呼ばれる）観点を提示する。これは、2つの測度間の optimal transport を、全長を最小化する測度空間内の曲線として記述するものである。動的な観点は optimal transport の代替的かつ直観的な解釈を提供し、流体力学との関連を引き出すだけでなく、2つの密度間の補間を行う際に低次元での OT 計算のための効率的な数値ツールをもたらす。この手法の欠点は、大規模な疎測度に対してスケールできず、二乗測地線コストを持つ正則な領域上の低次元でのみ機能すること（空間をグリッドで分割する必要があるため）である。

本章では、時刻 $t=0$ から1つの測度を出発し、時刻 $t=1$ にもう1つの測度に到達するという考え方に合わせて、$(\al,\be)$ の代わりに $(\al_0,\al_1)$ という記法を用いる。

%%%%%%%%%%%%%%%%%%%%%%%%%%%%%%%%%%%%%%%%%
\section{連続的定式化}

$\X=\Y=\RR^d$ かつ $c(x,y)=\norm{x-y}^2$ の場合、式~\eqref{eq-mk-generic}で定義された optimal transport 距離 $\Wass_2^2(\al,\be) = \MK_\c(\al,\be)$ は、2つの測度間の最小長パス $(\al_t)_{t=0}^1$ を探すことによって計算できる。
このパスは、各時刻に定義されるベクトル場 $\speed_t$ を用いて測度を移流することで記述される。ベクトル場 $\speed_t$ とパス $\alpha_t$ は質量保存式を満たさなければならず、その結果
\eql{\label{eq-cons-bb-ncvx}
	\pd{\al_t}{t} + \diverg(\al_t \speed_t) = 0
	\qandq
	\al_{t=0}=\al_0,
	\al_{t=1}=\al_1,
}
となる。ここで上の方程式は $\RR^d$ 上の分布（超関数）の意味で理解される。このようなベクトル場の微小長さは、測度 $\al_t$ に付随する $L^2$ ノルムを用いて測定され、次のように定義される。
\eq{\norm{\speed_t}_{L^2(\al_t)} = \left( \int_{\RR^\dim} \norm{\speed_t(x)}^2 \d\al_t(x) \right)^{1/2}.}
この定義から、\citet{benamou2000computational}によって最初に導入された $\Wass_2$ の以下の最小パス再定式化が得られる。
\eql{\label{eq-bb-non-cvx}
	\Wass_2^2(\al_0,\al_1) = \umin{ (\al_t,\speed_t)_t \text{ sat. \eqref{eq-cons-bb-ncvx}} }
		\int_0^1 \int_{\RR^d}   \norm{\speed_t(x)}^2 \d\al_t(x) \d t,
}
ここで $\al_t$ はスカラー値測度、$\speed_t$ はベクトル値測度である。図~\ref{fig-displacement}にこのような測度のパスの2つの例を示す。


\newcommand{\MyFigDispl}[1]{\includegraphics[width=.19\linewidth]{displacement/#1}}
\begin{figure}[h!]
\centering
\begin{tabular}{@{}c@{\hspace{1mm}}c@{\hspace{1mm}}c@{\hspace{1mm}}c@{\hspace{1mm}}c@{}}
\MyFigDispl{interp-dens-1} &
\MyFigDispl{interp-dens-2} &
\MyFigDispl{interp-dens-3} &
\MyFigDispl{interp-dens-4} &
\MyFigDispl{interp-dens-5} \\
\MyFigDispl{interp-points-1} &
\MyFigDispl{interp-points-2} &
\MyFigDispl{interp-points-3} &
\MyFigDispl{interp-points-4} &
\MyFigDispl{interp-points-5} \\
$\al_0$ & $\al_{1/4}$ & $\al_{1/2}$ & $\al_{3/4}$ & $\al_{1}$
\end{tabular}
\caption{\label{fig-displacement}
式~\eqref{eq-bb-non-cvx}を満たす変位補間 $\al_t$。
上段：Lebesgue 測度に関して密度を持つ2つの測度 $(\al_0,\al_1)$ に対するもの。
下段：同数の点を持つ2つの離散経験測度に対するもの。
}
\end{figure}


定式化~\eqref{eq-bb-non-cvx}は、制約~\eqref{eq-cons-bb-ncvx}が積 $\al_t \speed_t$ を含むため、変数 $(\al_t,\speed_t)_t$ に関して非凸である。ベクトル値測度（しばしば「運動量」と呼ばれる）
\eq{
	\Moment_t \eqdef \al_t \speed_t,
}
を導入することで、\citeauthor{benamou2000computational}は画期的な論文~\citeyearpar{benamou2000computational}において、次のように書くと変数 $(\al_t,\Moment_t)_t$ に関して凸となることを示した。
\eql{\label{eq-bb-convex}
	\Wass_2^2(\al_0,\al_1) = \umin{ (\al_t,\Moment_t)_t  \in \BBConstr(\al_0,\al_1)  }
		\int_0^1 \int_{\RR^d} \theta(\al_t(x),\Moment_t(x)) \d x \d t,
}
ここで制約集合は
\eql{\label{eq-mass-conservation}
	\BBConstr(\al_0,\al_1) \eqdef
	\enscond{ (\al_t,\Moment_t) }{
		\pd{\al_t}{t} + \diverg(\Moment_t) = 0,
		\al_{t=0}=\al_0,
		\al_{t=1}=\al_1
	},
}
と定義され、$\theta :  \rightarrow \RR^+ \cup \{+\infty\}$ は以下の下半連続凸関数である。
\eql{\label{eq-theta-dynamic}
	\foralls (a,b) \in \RR_+ \times \RR^\dim, \quad
	\th(a,b) = \choice{
		\frac{\norm{b}^2}{a} \qifq a>0, \\
		0 \qifq (a,b) = 0, \\
		+\infty \quad \text{otherwise}.
	}
}
この定義は複雑に見えるかもしれないが、$\al_t(x)=0$ のとき運動量 $\Moment_t(x)$ が消えることを課すために不可欠である。
また、式~\eqref{eq-bb-convex}は測度 $(\al_t,\Moment_t)$ が密度関数であるかのように非形式的に書かれているが、$\th$ は1次同次関数であり（したがって測度が Lebesgue 測度に関する密度を持たない場合でも定義され）、密度から関数への拡張が一意に行えるため、これは許容される。

%%%%%%%%%%%%%%%%%%%
\begin{rem11}{McCannの補間との関係}{rem-displacement}
最適 Monge 写像 $\T : \RR^\dim \rightarrow \RR^\dim$ が存在し（式~\eqref{eq-brenier-map}を参照）、$\T_\sharp \al_0=\al_1$ を満たす場合、$\al_t$ は McCann の補間
\eql{\label{eq-mc-cann-interp}
	\al_t = ((1-t) \Id + t \T)_{\sharp} \al_0.
}
に等しい。1次元の場合、注意~\ref{rem-1d-ot-generic}を用いると、この補間は以下の関係式により計算できる。
\eql{\label{eq-displacement-1d-cumul}
	\cumul{\al_t}^{-1} = (1-t)\cumul{\al_0}^{-1} + t \cumul{\al_1}^{-1};
}
図~\ref{fig-1d-ot}を参照。Wasserstein 空間のリーマン幾何学に関する詳細なレビューについては\citet{gangbo1996geometry}を参照されたい。
最適カップリング $\pi$ が存在するが、それが必ずしも Monge 写像の台上にない場合には、この補間を次のように計算できる。
\eql{\label{eq-displ-discr}
	\al_t = P_{t\sharp} \pi
	\qwhereq
	P_t : (x,y) \in \RR^\dim \times \RR^\dim \mapsto (1-t) x + t y.
}
例えば、離散的な設定~\eqref{eq-pair-discr}において、$\P$ を~\eqref{eq-mk-discr}の解とすると、補間は次のように定義される。
\eql{\label{eq-mccann-discrete}
	\al_t = \sum_{i,j} \P_{i,j} \de_{ (1-t) x_i + t y_j }.
}
このような補間は典型的に $n+m-1$ 個の点の上に台を持ち、これは $\P$ の非零要素数の最大値である。
図~\ref{fig-displacement-emp-weight}に離散測度の変位補間の2つの例を示す。
この構成は、$P_t$ を測地線パスに沿った補間に置き換えることで、測地的空間 $\Xx$ に一般化できる。
McCann の補間は、例えばコンピュータグラフィックスにおける色、形状、照明の補間~\citep{Bonneel-displacement}など、多くの応用がある。
\end{rem11}
%%%%%%%%%%%%%%%%%%%



\newcommand{\MyFigDisplEmp}[1]{\includegraphics[width=.16\linewidth]{interpolation-discrete/empirical/interp-#1}}
\newcommand{\MyFigDisplWei}[1]{\includegraphics[width=.16\linewidth]{interpolation-discrete/weighted/interp-#1}}
\begin{figure}[h!]
\centering
\begin{tabular}{@{}c@{}c@{}c@{}c@{}c@{}c@{}}
\MyFigDisplEmp{1}&
\MyFigDisplEmp{2}&
\MyFigDisplEmp{3}&
\MyFigDisplEmp{4}&
\MyFigDisplEmp{5}&
\MyFigDisplEmp{6}\\
\MyFigDisplWei{1}&
\MyFigDisplWei{2}&
\MyFigDisplWei{3}&
\MyFigDisplWei{4}&
\MyFigDisplWei{5}&
\MyFigDisplWei{6}\\
$\al_0$ & $\al_{1/5}$ & $\al_{2/5}$ & $\al_{3/5}$ & $\al_{4/5}$ & $\al_{1}$
\end{tabular}
\caption{\label{fig-displacement-emp-weight}
離散測度の変位補間~\eqref{eq-displ-discr}の比較。
上段：点群（同数の点を持つ経験測度 $(\al_0,\al_1)$）。
下段：同様であるが重みが異なる場合。$0<t<1$ において、上段の例は $N$ 個の点を持つ経験測度の補間 $\mu_t$ に対応し、下段は $2N-1$ 個の点の上に台を持つ測度を定義する。
}
\end{figure}

%%%%%%%%%%%%%%%%%%%%%%%%%%%%%%%%%%%%%%%%%
\section{一様スタガード格子上の離散化}

簡単のため、次元 $\dim=2$ の場合に数値スキームを説明する。より高次元への拡張は直截的である。ここでは\citet{FPapPeyOud13}によって導入された離散化法に従う。この方法は、流体力学で一般的に使用されるスタガード格子技法に着想を得ている。
時間を $t_k = k/T \in [0,1]$ と離散化し、空間は点 $x_i = (i_1/n_1,i_2/n_2) \in X = [0,1]^2$ で一様に離散化されるものとする。
スタガード格子表現を用い、$\al_t$ は時間方向の半格子点に関連付けられた $\a \in \RR^{(T+1) \times n_1 \times n_2}$ によって表現される。一方、$\Moment$ は $\MomentD=(\MomentD_1,\MomentD_2)$ によって表現され、$\MomentD_1 \in \RR^{T \times (n_1+1) \times n_2}$ と $\MomentD_1 \in \RR^{T \times n_1 \times (n_2+1)}$ は各空間方向の半格子点に格納される。
この表現を用いると、$(k,i_1,i_2) \in \range{1,T} \times \range{1,n_1} \times \range{1,n_2}$ に対して、時間微分は
\eq{
	(\partial_t \a)_{k,i} \eqdef \a_{k+1,i}-\a_{k,i}
}
と計算され、空間発散は
\eql{\label{eq-space-diverg-disc}
	\diverg(\MomentD)_{k,i} \eqdef
			\MomentD_{k,i_1+1,i_2}^1-\MomentD_{k,i_1,i_2}^1
			+
			\MomentD_{k,i_1,i_2+1}^2-\MomentD_{k,i_1,i_2}^2,
}
と計算される。これらはいずれも格子点上で定義され、$\RR^{T \times n_1 \times n_2}$ の配列を形成する。

最適化すべき汎関数を評価するために、半格子点から格子点への補間作用素が必要となる。すべての $(k,i_1,i_2) \in \range{1,T} \times \range{1,n_1} \times \range{1,n_2}$ に対して、
\eq{
	\InterpBB_a( \a )_{k,i} \eqdef \InterpBB(\a_{k+1,i},\a_{k,i}),
}
\eq{
	\InterpBB_\Moment( \MomentD )_{k,i} \eqdef (
		\InterpBB(\MomentD_{k,i_1+1,i_2}^1,\MomentD_{k,i_1,i_2}^1),
		\InterpBB(\MomentD_{k,i_1,i_2+1}^2,\MomentD_{k,i_1,i_2}^2)
	 ).
}
最も単純な選択は線形作用素 $\InterpBB(r,s)=\frac{r+s}{2}$ を用いることであり、以下ではこれを考える。式~\eqref{eq-bb-convex}の離散的対応物は次のようになる。
\eql{\label{eq-bb-discrete}
	\umin{(\a,\MomentD) \in \BBConstrD(\a_0,\a_1)}
		\Theta(\InterpBB_a(\a),\InterpBB_\Moment( \MomentD )),
}
\eq{
		\qwhereq
		\Theta(\tilde\a,\tilde\MomentD)
		\eqdef \sum_{k=1}^{T} \sum_{i_1=1}^{n_1} \sum_{i_2=1}^{n_2}
		\th( \tilde\a_{k,i}, \tilde\MomentD_{k,i} ),
}
ここで制約は
\eq{
	\BBConstrD(\a_0,\a_1) \eqdef
	\enscond{ (\a, \MomentD)
	}{
		\partial_t \a + \diverg(\MomentD) = 0,
		(\a_{0,\cdot},\a_{T,\cdot}) = (\a_0,\a_1)
	},
}
と読み替えられる。ここで $\a \in \RR^{(T+1) \times n_1 \times n_2}$、$\MomentD=(\MomentD_1,\MomentD_2)$ で $\MomentD_1 \in \RR^{T \times (n_1+1) \times n_2}$、$\MomentD_2 \in \RR^{T \times n_1 \times (n_2+1)}$ である。
図~\ref{fig-dynamic-maze}に、この離散化スキームを用いて近似された発展 $(\al_t)_t$ の例を示す。

%%%%
\begin{rem}[グラフ上の動的定式化]
$\X$ がグラフであり $c(x,y)=d_\X(x,y)^2$ が二乗測地線距離である場合、一様格子離散化~\eqref{eq-space-diverg-disc}の代わりにグラフ構造に付随する離散発散を用いた忠実な離散化手法を導出することが可能である。対応する動的 Wasserstein 距離に対して熱方程式が勾配流構造を持つことを保証するために（勾配流の詳細については第~\ref{sec-grad-flows}節を参照）、\citet{Maas2011}および後に\citet{MielkeCVPDE}は対数平均 $\InterpBB(r,s)$ を用いることを提案した（\citep{solomon2016continuous,ChowHuangLiZhou2012,chow2017entropy,chow2017discrete}も参照）。
\end{rem}
%%%%

%%%%%%%%%%%%%%%%%%%%%%%%%%%%%%%%%%%%%%%%%
\section{近接法ソルバー}
\label{sec-prox-solvers}

離散化された動的 OT 問題~\eqref{eq-bb-discrete}は、アフィン制約の下で非滑らかな最適化問題を最小化する必要があるため、解くことが困難である。実際、関数 $\th$ は凸であるが、消滅する質量 $\a_{k,i}$ を持つ測度に対しては非滑らかである。コンパクトな台を持つ2つの入力測度 $(\a_0,\a_1)$ 間の補間において、補間された測度 $(\a_{k})_{k=1}^T$ の質量もまた消滅することが典型的に予想され、最適化過程の困難な部分はまさにこの台の発展を追跡することにある。特に、標準的な滑らかな最適化手法を使用することはできない。

式~\eqref{eq-bb-discrete}を二次錐計画に書き直す方法がいくつかある。双対問題を考えるか、あるいは単に汎関数 $\th( \a_{k,i}, \MomentD_{k,i} )$ を制約付きの線形関数に置き換えることによる。
\eq{
	\Theta(\tilde\a,\tilde\MomentD) = \umin{\tilde \zD} \enscond{
		\sum_{k,i} \tilde \zD_{k,i}
	}{
		\foralls (k,i), (\zD_{k,i},\tilde\a_{k,i},\tilde\MomentD_{i,j}) \in \Ll
	},
}
これには追加変数 $\tilde \zD$ の導入が必要となる。
ここで $\Ll \eqdef \ensconds{ (z,a,J) \in \RR \times \RR^+ \times \RR^\dim }{ \norm{J}^2 \leq z a }$ は回転 Lorentz 二次錐である。
この追加変数を用いることで、二次錐計画に対する標準的な内点法ソルバー~\citep{nesterov1994interior}を使用して離散化問題を解くことが可能となる。これらのソルバーは高速な収束率を持ち、高精度の解を計算できる。しかしながら、各反復はコストが高く、離散化点数にスケールする次元の線形系の求解を必要とする。したがって、イメージング応用で遭遇する大規模多次元問題には適用できない。

これらの高精度ソルバーの代替となるのが、低精度の一次法であり、式~\eqref{eq-bb-discrete}のような非滑らかだが高度に構造化された問題に適している。このクラスのソルバーは新しいものではないが、最近、イメージングと機械学習の分野で再び注目されている。これらの応用では数値精度が最優先の目標ではないためである。
これらのソルバーの詳細な説明と大規模応用への使用については、例えばモノグラフ~\citep{BauschkeCombettes11}を参照されたい。
ここでは特定のソルバーに焦点を当てるが、もちろん他の多くのソルバーも使用可能であり、動的 OT に対するこれらのアプローチのいくつかの研究については\citep{FPapPeyOud13}を参照されたい。動的 OT に対する一次スキームの使用は、最初に\citet{benamou2000computational}によって提案されたことに注意する。

DR アルゴリズム~\citep{Lions-Mercier-DR}は、以下の形の非滑らかな構造化問題を解くために特別に設計されている。
\eql{\label{eq-splitting-opt}
	\umin{x \in \Hh} F(x) + G(x),
}
ここで $\Hh$ はあるユークリッド空間であり、$F, G : \Hh \rightarrow \RR \cup \{+\infty\}$ は2つの閉凸関数であって、パラメータ $\tau>0$ に対するいわゆる近接作用素
\eql{\label{eq-prox-eucl}
	\foralls x \in \Hh, \quad
	\Prox_{\tau F}(x) \eqdef \uargmin{x' \in \Hh} \frac{1}{2}\norm{x-x'}^2 + \tau F(x)
}
を「容易に」（\eg 閉じた形で、あるいは急速に収束するスキームを用いて）計算できるものである。これはユークリッド計量に対する近接写像に対応し、この定義は $\norm{x-x'}^2$ の代わりにより一般的な Bregman ダイバージェンスに拡張できることに注意する。$\KLD$ ダイバージェンスを用いた例については式~\eqref{eq-prox-kl}を参照。
DR アルゴリズムの反復は、初期化 $(x^{(0)}, w^{(0)})\in \Hh^2$ を用いて列 $(\it{x}, \it{w}) \in \Hh^2$ を以下のように定義する。
\begin{equation}\label{eq-dr-iters}
\begin{aligned}
	\itt{w} &\eqdef  \it{w} + \al (\Prox_{\ga F} (2\it{x}-\it{w})-\it{x}),\\
	\itt{x} &\eqdef  \Prox_{\ga G}(\itt{w}).
\end{aligned}
\end{equation}
$0 < \al < 2$ かつ $\ga > 0$ であれば、$\it{x} \rightarrow z^\star$ が示され、ここで $z^\star$ は式~\eqref{eq-splitting-opt}の解である。詳細については\citep{Combettes2007}を参照。
このアルゴリズムは、もう1つの広く使われている手法である交互方向乗数法~\citep{GabayMercier,GlowinskiMarroco}（レビューとして\citep{BoydADMM}も参照）と密接に関連しており、DR を双対問題に適用することで得られる。2つの手法間の同値性については\citep{FPapPeyOud13}を参照されたい。この同値性は最初に\citep{Eckstein1992}によって示された。

問題~\eqref{eq-bb-discrete}を形式~\eqref{eq-splitting-opt}に書き直す方法は多数あり、これらのアプローチの詳細な説明については\citep{FPapPeyOud13}を参照されたい。これを達成する簡単な方法は、$x=(\a,\MomentD,\tilde\a,\tilde\MomentD)$ と設定し、
\eq{
	F(x) \eqdef \Theta(\tilde\a,\tilde\MomentD) + \iota_{\BBConstrD(\a_0,\a_1)}(\a,\MomentD)
	\qandq
	G(x) = \iota_{\Dd}(\a,\MomentD,\tilde\a,\tilde\MomentD),
}
\eq{
	\qwhereq \Dd \eqdef \enscond{(\a,\MomentD,\tilde\a,\tilde\MomentD)}{
		\tilde\a=\InterpBB_a(\a), \tilde\MomentD=\InterpBB_\Moment( \MomentD )
	}.
}
とすることである。
これら2つの関数の近接作用素は効率的に計算できる。
実際、
\eq{
	\Prox_{\tau F}(x) = ( \Prox_{\tau\Theta}(\tilde\a,\tilde\MomentD), \Proj_{\BBConstrD(\a_0,\a_1)}(\a,\MomentD) ).
}
が成り立つ。近接作用素 $\Prox_{\tau\Theta}$ は各格子位置で3次多項式方程式を解くことで計算される。アフィン制約 $\BBConstrD(\a_0,\a_1)$ への直交射影は Poisson 方程式の求解を含み、高速 Fourier 変換を用いて $O(N\log(N))$ 演算で達成できる。ここで $N=Tn_1 n_2$ は格子点数である。
最後に、近接作用素 $\Prox_{\tau G}$ は線形射影子であり、小さな線形系の逆行列計算を必要とする。
これらの計算の詳細については\citet{FPapPeyOud13}を参照されたい。
図~\ref{fig-dynamic-maze}に、この手法を用いて複雑な平面領域内で動的補間を計算した例を示す。
動的 OT に対するこのクラスの近接法は、平均場ゲーム~\citep{benamou2015augmented}などの関連問題を解くためにも使用されている。

\begin{figure}[h!]
\centering
\includegraphics[width=\linewidth]{dynamic/dynamic-maze}
\caption{\label{fig-dynamic-maze}
近接分裂スキームを用いて計算された動的 OT の解 $\al_t$。
}
\end{figure}



%%%%%%%%%%%%%%%%%%%%%%%%%%%%%%%%%%%%%%%%%
\section{動的非均衡 OT}
\label{dynamic-unbalanced}

異なる質量 $\al_0(\Xx) \neq \al_1(\Xx)$ を持つ入力測度を対応付けるため（いわゆる「非均衡」設定、\citet{benamou2003numerical}によって導入された用語）、また局所的な質量変動に対処するため、いくつかの正規化や緩和が提案されている。特に固定周辺制約を緩和することによるものがある。第~\ref{sec-unbalanced}節を参照。
一般的な方法論は、連続の方程式~\eqref{eq-mass-conservation}にソース項 $s_t(x)$ を導入することからなる。したがって
\eq{
	\bar\BBConstr(\al_0,\al_1) \eqdef
	\enscond{ (\al_t,\Moment_t,s_t) }{
		\pd{\al_t}{t} + \diverg(\Moment_t) = s_t,
		\al_{t=0}=\al_0,
		\al_{t=1}=\al_1
	}.
}
を考える。重要な問題は、このソース項に関連するコストをどのように測定し、元の動的定式化~\eqref{eq-bb-convex}に導入するかである。文献にはいくつかの提案がある。例えば $L^2$ コスト~\citet{piccoli2014generalized}を用いるものがある。
質量を「テレポート」させること（質量が無限の速度で移動し、それまで質量がなかった領域に突然現れること）を避けるために、関連するコストは無限大であるべきである。これは、$s_t$ が任意の測度であることを許し（\emph{例えば}1次同次コストを用いて）、$s_t$ を運動量 $\Moment_t$ と同じ方法で罰則化することで、単純な凸的方法で達成できることが判明する。
\eql{\label{eq-wfr-dynamic}
	\WFR^2(\al_0,\al_1) = \umin{ (\al_t,\Moment_t,s_t)_t  \in \bar\BBConstr(\al_0,\al_1)  }
		\Theta(\al,\Moment,s),
}
\eq{
	\qwhereq \Theta(\al,\Moment,s) \eqdef
		\int_0^1 \int_{\RR^d} \pa{
			\theta(\al_t(x),\Moment_t(x)) + \tau \theta(\al_t(x),s_t(x))
		}\d x \d t,
}
ここで $\th$ は式~\eqref{eq-theta-dynamic}で導入された凸1次同次関数であり、$\tau$ は質量輸送と質量生成・消滅のトレードオフを制御する重みである。この定式化は複数の著者によって独立に提案された~\citep{LieroMielkeSavareShort,2017-chizat-focm,kondratyev2015}。この「動的」定式化には「静的」な対応物がある。注意~\ref{rem-wfr-static}を参照。
凸最適化問題~\eqref{eq-wfr-dynamic}は、第~\ref{sec-prox-solvers}節で詳述した方法と同様の手法で解くことができる。
図~\ref{fig-unbalanced-dynamic}に、いくつかの非均衡 OT の動的補間の比較を示す。
この動的定式化は、形状レジストレーションのための「メタモルフォシス」モデル~\citep{Metamorphosis2005}に類似しており、より正確な関連は\citep{maas2015generalized,maas2016generalized}で詳述されている。

$\tau \rightarrow 0$ とし、$\al_0(\Xx) = \al_1(\Xx)$ であれば、古典的な OT 問題が回復される。すなわち $\WFR(\al_0,\al_1) \rightarrow \Wass(\al_0,\al_1)$ となる。一方、$\tau \rightarrow +\infty$ では、この距離は密度上の Hellinger 計量に近づく。
\begin{align*}
	\frac{1}{\tau}\WFR(\al_0,\al_1)^2\; \overset{\tau \rightarrow +\infty} \longrightarrow
	&\int_{\Xx} | \sqrt{\density{\al_0}(x)} - \sqrt{\density{\al_1}(x)} |^2 \d x \\
	=& \int_{\Xx} | 1 - \sqrt{ \frac{\d\al_1}{\d\al_0}(x) } |^2 \d \al_0(x).
\end{align*}

\begin{figure}[h!]
\centering
\includegraphics[width=\linewidth]{unbalanced/unbalanced-dynamic}
\caption{\label{fig-unbalanced-dynamic}
Hellinger（第1行）、Wasserstein（第2行）、部分 optimal transport（第3行）、および Wasserstein--Fisher--Rao（第4行）の動的補間の比較。
}
\end{figure}




%%%%%%%%%%%%%%%%%%%%%%%%%%%%%%%%%%%%%%%%%
\section{より一般的な移動性汎関数}
\label{eq-generic-modibility}

動的定式化~\eqref{eq-bb-convex}を、式~\eqref{eq-theta-dynamic}で定義されたものに代えて他の「移動性関数」$\th$ を考えることによって一般化することが可能である。この移動性汎関数の可能な選択は\citet{dolbeault2009new}で提案されている。
\eql{\label{eq?generic-mobility}
	\foralls (a,b) \in \RR_+ \times \RR^\dim, \quad
	\th(a,b) = a^{s-p} \norm{b}^{p},
}
ここでパラメータは $\th$ が凸となるために $p \geq 1$ かつ $s \in [1,p]$ を満たす必要がある。
$1<s\leq p$ の場合には $\th$ が無限遠で線形増大を持たないため注意が必要であり、式~\eqref{eq-bb-convex}の解は Lebesgue 測度に関する密度を持つことが制約される。

$s=1$ の場合は古典的な OT 問題に対応し、式~\eqref{eq-bb-convex}の最適値は $\Wass_p(\al,\be)$ を定義する。この場合、$\th$ は1次同次であり、式~\eqref{eq-bb-convex}の解は任意の測度であり得る。$(s=1,p=2)$ の場合が、$\Wass_2$ を定義するために式~\eqref{eq-bb-convex}で考えられた初期設定である。

極限の場合 $s=p$ も興味深く、双対 Sobolev ノルム $W^{-1,p}$ に対応する。式~\eqref{eq-bb-convex}の値は次のようになる。
\eq{
	\norm{\al-\be}_{W^{-1,p}(\RR^\dim)}^p =
	\umin{\f} \enscond{ \int_{\RR^\dim} \f  \d(\al-\be)}{  \int_{\RR^\dim} \norm{ \nabla f(x) }^q \d x \leq 1 }
}
ただし $1/q+1/p=1$ である。
極限 $(p=s,q) \rightarrow (1,\infty)$ において、$\Wass_1$ ノルムが回復される。
$s=p=2$ の場合は、式~\eqref{eq-dual-sobolev-div}で定義された Sobolev $H^{-1}(\RR^\dim)$ Hilbert ノルムに対応する。


%%%%%%%%%%%%%%%%%%%%%%%%%%%%%%%%%%%%%%%%%
\section{パス空間上の動的定式化}
\label{sec-entropic-dynamic}

古典的および正則化されたエントロピー正則化（第~\ref{c-entropic}章を参照）による OT の定式化には自然な動的定式化が存在し、これは空間 $\Xx$ 上のすべての可能なパス $\ga : [0,1] \rightarrow \Xx$（\ie 曲線）の空間 $\bar\Xx$ 上の抽象的な最適化問題の研究に基づいている。
簡単のため $\Xx=\RR^\dim$ と仮定するが、これは測地的空間やグラフなどのより一般的な空間にも拡張される。
非形式的に、時刻 $t=0,1$ における2つの入力測度 $\al_0,\al_1$ の間の「粒子」の動力学は、確率分布 $\bar\pi \in \Mm_+^1(\bar\Xx)$ によって記述される。このような分布は、始点と終点の分布が $(\al_0,\al_1)$ と一致しなければならないという条件を満たす必要があり、これは押し出し写像を用いて形式的に
\eq{
	\bar\Couplings(\al_0,\al_1) \eqdef \enscond{ \bar\pi \in \Mm_+^1(\bar\Xx) }{ \bar P_{0\sharp} \bar\pi = \al_0, \bar P_{1\sharp} \bar\pi = \al_1  },
}
と書かれる。ここで任意のパス $\ga \in \bar\Xx$ に対して $P_0(\ga) = \ga(0), P_1(\ga) = \ga(1)$ である。

%%%%
\paragraph{パス空間上の OT}

古典的 OT~\eqref{eq-mk-generic}のパス空間上の動的版は次のように読める。
\eql{\label{eq-ot-pathsspace}
	\Wass_2(\al_0,\al_1)^2 =
	\umin{ \bar\pi \in \bar\Couplings(\al_0,\al_1) } \int_{\bar\Xx} \Ll(\ga)^2 \d\bar\pi(\ga),
}
ここで $\Ll(\ga) = \int_0^1 |\ga'(s)|^2 \d s$ はパス $s \in [0,1] \mapsto \ga(s) \in \Xx$ の運動エネルギーである。それぞれ式~\eqref{eq-ot-pathsspace}と式~\eqref{eq-mk-generic}を解く最適カップリング $\pi^\star$ と $\bar\pi^\star$ の間の関係は、$\bar\pi^\star$ が $\pi^\star$ によって規定される割合で点の対を結ぶ測地線にのみ質量を与えるというものである。離散測度の特別な場合には、次のようになる。
\eq{
	\pi^\star = \sum_{i,j} \P_{i,j} \de_{(x_i,y_j)}
	\qandq
	\bar\pi^\star = \sum_{i,j} \P_{i,j} \de_{\ga_{x_i,y_j}},
}
ここで $\ga_{x_i,y_j}$ は $x_i$ と $y_j$ の間の測地線である。さらに、$\bar\pi^\star$ に従って引かれた $\ga$ の時刻 $t$ における曲線点 $\ga(t)$ の分布によって定義される測度、\ie
\eql{\label{eq-interpolation-spacepaths}
	 t \in [0,1] \mapsto \al_t \eqdef P_{t\sharp} \bar\pi^\star
	 \qwhereq
	 P_t(\ga) = \ga(t) \in \Xx,
}
は動的定式化~\eqref{eq-bb-convex}の解、\ie 変位補間である。
離散的な場合には、式~\eqref{eq-mccann-discrete}が回復される。

%%%%
\paragraph{パス空間上のエントロピー正則化 OT}

ここでは、第~\ref{c-entropic}章で定義されたエントロピー正則化 OT のパス空間を用いた再解釈に移る。
式~\eqref{eq-entropic-generic-proj}と同様に、これはパス空間上の参照測度 $\bar\Kk$ に対する Kullback--Leibler 射影を用いて定義される。ここで $\bar\Kk$ は、初期時刻と最終時刻に一様分布を持つ可逆 Brown 運動（Wiener 過程）の分布である。
\eql{\label{eq-shrodinger-paths}
	\umin{ \bar\pi \in \bar\Couplings(\al_0,\al_1) } \KL(\bar\pi|\bar\Kk).
}
この問題の概要と\citet{Schroedinger31}の研究の歴史的説明については、\citet{LeonardSchroedinger}によるレビュー論文を参照されたい。
式~\eqref{eq-shrodinger-paths}の（一意の）解 $\bar\pi_\varepsilon^\star$ は $\varepsilon \rightarrow 0$ のとき式~\eqref{eq-ot-pathsspace}の解に収束することが示される。さらに、この解は Brown 橋 $\bar\ga_{x,y}^\varepsilon \in \bar\Xx$（ぼやけた測地線に類似し、$\varepsilon \rightarrow 0$ のとき $\de_{\ga_{x,y}}$ に収束する）を用いた静的エントロピー正則化 OT 問題~\eqref{eq-entropic-generic}の解と関連している。離散的な設定では、次のようになる。
\eql{\label{eq-brownian-bridge-interp}
	\pi_\varepsilon^\star = \sum_{i,j} \P_{\varepsilon,i,j}^\star \de_{(x_i,y_j)}
	\qandq
	\bar\pi_\varepsilon^\star = \sum_{i,j} \P_{\varepsilon,i,j}^\star \bar\ga_{x_i,y_j}^\varepsilon,
}
ここで $\P_{\varepsilon,i,j}^\star$ は Sinkhorn アルゴリズムを用いて計算できる。
式~\eqref{eq-interpolation-spacepaths}と同様に、エントロピー正則化補間を次のように定義できる。
\eq{
	\al_{\varepsilon,t} \eqdef \P_{t\sharp} \bar\pi_\varepsilon^\star.
}
Brown 橋に沿った時刻 $t$ の位置の法則 $\P_{t\sharp} \bar\ga_{x,y}^\varepsilon$ は $\ga_{x,y}(t)$ を中心とする分散 $t(1-t)\varepsilon^2$ のガウス分布 $\Gg_{t(1-t)\varepsilon^2}(\cdot-\ga_{x,y}(t))$ であるから、$\al_{\varepsilon,t}$ は移動する Dirac 質量の集合のガウスぼかしであることが導かれる。
\eq{
	\al_{\varepsilon,t} =
	\sum_{i,j} \P_{\varepsilon,i,j}^\star \Gg_{t(1-t)\varepsilon^2}( \cdot - \ga_{x_i,y_j}(t) ).
}
結果として得られる Brown 橋の混合を図~\ref{fig-schrodinger-dynamic}に示す。


\newcommand{\MyFigSchrDyn}[1]{\includegraphics[width=.245\linewidth,trim=50 45 38 35,clip]{schrodinger-dynamic/schrodinger-dynamic-eps#1}}
%%% FIG %%%
\begin{figure}[h!]\centering
\centering
\begin{tabular}{@{}c@{}c@{}c@{}c@{}}
\MyFigSchrDyn{5}&
\MyFigSchrDyn{50}&
\MyFigSchrDyn{200}&
\MyFigSchrDyn{1000}\\
$\varepsilon=0$ & $\varepsilon=.05$ & $\varepsilon=0.2$ & $\varepsilon=1$
\end{tabular}
%%%
\caption{%%%
Schr\"odinger エントロピー正則化補間~\eqref{eq-brownian-bridge-interp}に関連するパス空間上の Brown 橋パスのサンプル。青は $t=0$ に、赤は $t=1$ に対応する。
} \label{fig-schrodinger-dynamic}
\end{figure}
%%% FIG %%%

このエントロピー正則化補間 $(\al_t)_t$ を記述するもう1つの方法は、Benamou--Brenier 動的定式化~\eqref{eq-bb-non-cvx}の正則化を用いるものである。すなわち
\eql{\label{eq-bb-non-cvx-entropy}
	\umin{ (\al_t,\speed_t)_t \text{ sat. \eqref{eq-cons-bb-ncvx}} }
		\int_0^1 \int_{\RR^d}  \pa{
			 \norm{\speed_t(x)}^2  + \frac{\varepsilon}{4} \norm{ \nabla \log(\al_t)(x) }^2
			 }
			 \d\al_t(x) \d t;
}
\citep{gentil2015analogy,chen2016relation}を参照。
